\documentclass[reqno]{amsart}
\usepackage[T1]{fontenc}
\usepackage{amsmath,amssymb,amsthm,mathtools}
\usepackage[shortlabels]{enumitem}
\usepackage{xr-hyper}
\usepackage[hidelinks]{hyperref}
\externaldocument[ft-]{20261124T213000--schwartz-space-cover-exercises__workbook}
\externaldocument[in-]{20261124T213100--schwartz-space-cover-inversion__workbook}

\providecommand{\vol}{\operatorname{vol}}
\providecommand{\qr}[2]{\left(\frac{#1}{#2}\right)}
\newcommand{\SL}{\mathrm{SL}_2}
\newcommand{\tSL}{\widetilde{\mathrm{SL}}_2}
\newcommand{\Fc}{\mathcal{F}}
\newcommand{\Dc}{\mathcal{D}}
\newcommand{\Sc}{\mathcal{S}}
\newcommand{\Wc}{\mathcal{W}}
\newcommand{\Ab}{\mathbb{A}}
\newcommand{\oo}{\mathfrak{o}}
\newcommand{\pp}{\mathfrak{p}}

% Chapter number; \ref{chap} gives it.
\newcounter{chap}
\renewcommand{\thechap}{3}
\renewcommand{\thesection}{\thechap.\arabic{section}}
\numberwithin{equation}{section}
\theoremstyle{definition}
\newtheorem{problem}{Problem}[section]
\theoremstyle{plain}
\newtheorem{theorem}{Theorem}[section]

\newenvironment{hints}{\begin{description}[font=\normalfont\bfseries]}{\end{description}}
\newenvironment{exercises}{\subsection*{Exercises}\begin{enumerate}[1.]}{\end{enumerate}}

\setlength{\emergencystretch}{3em}
\title[A Poisson formula for the $n$-fold cover]{Leading terms and a Poisson
  formula for the $n$-fold cover of $\mathrm{SL}_2$}
\author{}
\date{November 24, 2026}

\begin{document}
\maketitle
\refstepcounter{chap}\label{chap}

For $n = 1$ and a global field $F$, Poisson summation on $F^2 \subset \Ab^2$ reads $\sum_{v \neq 0} \Phi(v) + \Phi(0) = \sum_{v \neq 0} \Fc \Phi(v) + \Fc \Phi(0)$.  For the $n$-fold cover of $\SL$, whose rational points carry Kubota's $n$th power residue symbols, what replaces the values at the origin?  This sheet defines a leading term at the origin, sums over the rational points of the plane, and proves a Poisson formula with a boundary term, following \cite{N}.

\paragraph{Prerequisites.}
The sheets \emph{From Poisson summation on the plane to a Fourier transform for the $n$-fold cover of $\mathrm{SL}_2$} (Chapter~\ref{ft-chap}) and \emph{The Schwartz space of the $n$-fold cover of $\mathrm{SL}_2$: symmetries, the basic vector and inversion} (Chapter~\ref{in-chap}), whose notation we keep.  Poisson summation on the adeles \cite[\S4.2]{Tate}, and Hilbert's reciprocity law \cite[Ch.~VI, Theorem~8.1]{Neu}.  For \S\ref{sec:poisson} only, reduction theory for $\SL$ over a global field \cite[I.2.1]{MW}.


%======================================================================
\section{Leading terms at the origin}\label{sec:germs}

For $\Psi \in \Sc$ and $v$ in either chart, the function $\mu \mapsto \lvert \mu \rvert^{n-1} \Psi(\mu^n v)$ extends to a Schwartz--Bruhat function on $F$, uniformly for $v$ in compact sets: for $\Psi \in \Dc$ by definition, and for $\Psi \in \Fc \Dc$ by Problem~\ref{ft-pr:ft-root} and Theorem~\ref{ft-th:lines} \cite[Proposition~8.1]{N}.  Its value at $\mu = 0$ defines the \emph{leading term} $b \Psi$, a genuine function.  It vanishes on $\Dc$.

\begin{problem}\label{pr:leading-transform}
  For $\Phi \in \Dc$, $(b \Fc \Phi)_X(v) = \lvert n \rvert^{-1/2} \int_F R_X \Phi(v, \tau) \, d \tau$.  For $n = 1$ this is $\int_{F^2} \Phi = \Fc \Phi(0)$.
\end{problem}

\begin{problem}\label{pr:leading-homog}
  $(b \Psi)(\mu^n v) = \lvert \mu \rvert^{1-n} (b \Psi)(v)$, and $b(\Psi(a \, \cdot)) = \lvert a \rvert^{-1 + 1/n} b \Psi$ for $a \in F^{\times n}$.
\end{problem}

Also $b(g \Psi) = g \, b \Psi$ for $g \in \tSL(F)$, and in the tame case $b \Phi^0 = \Phi^0$ on $\lVert v \rVert \leq 1$.

\begin{problem}\label{pr:finite-germs}
  If $F$ is non-archimedean, each $\Psi \in \Sc$ vanishes outside a compact subset of $F^2$ and agrees with $b \Psi$ near the origin, and the kernel of $b$ on $\Sc$ is $\Dc$.
\end{problem}

The map $\mu \mapsto \mu^n$ is $n$-to-one from $F^\times$ onto $F^{\times n}$ and multiplies $d^\times \mu$ by $\lvert n \rvert$, so $\int_{F^{\times n}} \phi(a) \, d^\times a = \frac{\lvert n \rvert}{n} \int_{F^\times} \phi(\mu^n) \, d^\times \mu$.  For $\Phi \in \Dc$, let $P \Phi := \int_{F^{\times n}} \Phi(a \, \cdot) \lvert a \rvert^{1 + 1/n} \, d^\times a$; it satisfies $(P \Phi)(a \, \cdot) = \lvert a \rvert^{-1-1/n} P \Phi$ for $a \in F^{\times n}$.

\begin{problem}\label{pr:germ-average}
  For $\Phi \in \Dc$, $b \Fc \Phi = n \lvert n \rvert^{-3/2} A(P \Phi)$.
\end{problem}

At a finite place every genuine $f$ with $f(a v) = \lvert a \rvert^{-1-1/n} f(v)$ for $a \in F^{\times n}$ is of the form $P \Phi$, and we obtain \cite[Corollary~8.4]{N}:

\begin{theorem}\label{th:finite-schwartz}
  Let $F$ be non-archimedean.  Then $b$ identifies $\Sc / \Dc$ with the image under $A$ of the genuine functions $f$ with $f(a v) = \lvert a \rvert^{-1-1/n} f(v)$ for $a \in F^{\times n}$.  A genuine function lies in $\Sc$ if and only if it vanishes outside a compact subset of $F^2$ and agrees near the origin with such an $A f$.
\end{theorem}

We say that $\Phi \in \Dc$ satisfies the \emph{vanishing condition} if $P \Phi = 0$; then $b \Phi = b \Fc \Phi = 0$.  For $n = 1$, Theorem~\ref{th:finite-schwartz} says that $\Sc$ consists of the restrictions of $C_c^\infty(F^2)$, and $b$ is evaluation at the origin.

Globally, the Poisson formula should be a sum over $F^2 - \{0\}$ of the values of $\Phi = \otimes_\pp \Phi_\pp$.  At a rational point $(x, y)$ with $x y \neq 0$, the product of the $X$-components and that of the $Y$-components differ by $\prod_\pp \{x, y\}_{\pp, n}$, which is $1$ by Hilbert's reciprocity law.

%======================================================================
\section{The adelic Schwartz space and its theta sum}\label{sec:adelic}

Let $F$ be a global field containing a primitive $n$th root of unity, $\Ab$ its adele ring, $\psi$ a nontrivial unitary character of $\Ab / F$, and $S$ a nonempty finite set of places such that $(F_\pp, \psi_\pp, n)$ is good for $\pp \notin S$.  Hilbert's reciprocity law states that for $x, y \in F^\times$, $\{x, y\}_{\pp, n} = 1$ for almost all $\pp$ and $\prod_\pp \{x, y\}_{\pp, n} = 1$ \cite[Ch.~VI, Theorem~8.1]{Neu} (for number fields; \cite{N} also treats function fields).  Let $\Sc_\Ab := \bigotimes'_\pp \Sc_\pp$, restricted with respect to the basic vectors $\Phi_\pp^0$.

\begin{problem}\label{pr:global-values}
  For a pure tensor $\Phi$ and $v \in F^2 - \{0\}$, the products of the local $X$-components and of the local $Y$-components define a single value $\Phi(v)$.  The factors $\lvert n \rvert_\pp^{-1/2}$ in $\Fc \Phi := \otimes \Fc \Phi_\pp$ do not affect its values at rational points.
\end{problem}

\begin{problem}\label{pr:theta-S}
  Suppose that $\mathcal{O}_S$ is a principal ideal domain and $\Phi_\pp = \Phi_\pp^0$ for $\pp \notin S$.  Then
  \begin{equation*}
    \sum_{v \in F^2 - \{0\}} \Phi(v) = \sum_{\tau} \lvert \tau \rvert_S^{n-1} \sum_{(c, d)} \qr{d}{c}_n \Phi_S(\tau^n c, \tau^n d),
  \end{equation*}
  where $\tau$ runs over $(\mathcal{O}_S - \{0\}) / \mathcal{O}_S^\times$, $(c, d)$ over coprime pairs in $\mathcal{O}_S$, $\lvert \tau \rvert_S := \prod_{\pp \in S} \lvert \tau \rvert_\pp$, and $\Phi_S$ is the product of the $X$-components at $\pp \in S$ (of the $Y$-components when $c = 0$, in which case the symbol is omitted).
\end{problem}

For $F = \mathbb{Q}(\omega)$, $n = 3$ and $S = \{\infty, (1 - \omega)\}$, the inner sum is a sum of cubic residue symbols over coprime pairs, as in Kubota's Eisenstein series.  The sum $\Theta(\Psi) := \sum_{v \neq 0} \Psi(v)$ converges absolutely for every $\Psi \in \Sc_\Ab$; this follows from a count of rational points in adelic boxes, using \eqref{ft-eq:growth} and the product formula \cite[Lemma~B.1]{N}.

At a finite place $\pp \nmid n$, call an element of $\tSL(F_\pp)$ \emph{integral} if it is a product of operators $n(u)$, $n^-(u)$ with $u \in \oo_\pp$.  A \emph{lift} of $m \in \SL(\Ab)$ is a family $(g_\pp)$ with $\pi_\pp(g_\pp) = m_\pp$ and $g_\pp$ integral for almost all $\pp$.  Two lifts of the same $m$ act on $\Sc_\Ab$ by maps differing by a scalar in $\mu_n$, and lifts commute with $\Fc$.

\begin{problem}\label{pr:rational-invariance}
  For $u \in F$ and $v \in F^2 - \{0\}$, $(n(u) \Psi)(v) = \Psi(v \, \pi(n(u)))$ and $(n^-(u) \Psi)(v) = \Psi(v \, \pi(n^-(u)))$.
\end{problem}

Since these matrices generate $\SL(F)$, each $\gamma \in \SL(F)$ has a lift $\tilde{\gamma}$ with $\Theta(\tilde{\gamma} \Psi) = \Theta(\Psi)$ \cite[Lemma~B.5]{N}.

We must compare $\Theta(\Phi)$ with $\Theta(\Fc \Phi)$.  The usual proof of Poisson summation uses Fourier inversion on a compact quotient group, but genuine functions are not functions on a group.  Their difference, evaluated on the translates of $\Phi$, is however an automorphic function, and the first thing to compute for an automorphic function is its constant term.

%======================================================================
\section{The Poisson formula}\label{sec:poisson}

For a pure tensor $\Psi$ and $x, y \in F^\times$, let $m_\Psi(x) := \prod_\pp m_{\Psi_\pp}(x)$, $r_\Psi(y) := \prod_\pp r_{\Psi_\pp}(y)$, and $C(\Psi) := \sum_{y \in F^\times} r_\Psi(y) + \sum_{x \in F^\times} m_\Psi(x)$.

\begin{problem}\label{pr:constant-term}
  $\int_{\Ab / F} \Theta(n(u) \Psi) \, d u = C(\Psi)$.
\end{problem}

To compare $C(\Phi)$ with $C(\Fc \Phi)$ we apply Poisson summation on $\Ab$ \cite[\S4.2]{Tate} to the two relations of \S\ref{in-sec:inversion}, and then Poisson summation on the $n$th power classes.  The group $F_\pp^\times / F_\pp^{\times n}$ has order $n^2 / \lvert n \rvert_\pp$ \cite[Ch.~II, Corollary~5.8]{Neu} and is self-dual under the Hilbert symbol, and the image of $F^\times$ in $\Ab^\times / \Ab^{\times n}$ is its own orthogonal complement \cite[Ch.~XIII, \S5, Proposition~8]{Weil}.  Let $\Delta := F^\times / F^{\times n}$.

\begin{problem}\label{pr:isotropy}
  For each $\pp$ let $\beta_\pp$ be a function on $F_\pp^\times / F_\pp^{\times n}$, equal to the indicator of the image of $\oo_\pp^\times$ for almost all $\pp$, and let $\tilde{\beta}_\pp(y) := \frac{1}{n} \sum_c \{c, y\}_{\pp, n} \beta_\pp(c)$.  Then $\sum_{x \in \Delta} \prod_\pp \beta_\pp(x) = \sum_{y \in \Delta} \prod_\pp \tilde{\beta}_\pp(y)$.
\end{problem}

For $y' \in F^\times$, let $\nu := \otimes_\pp \nu_\pp$ with $\nu_\pp(\tau) := R_Y \Phi_\pp((0, y'), \tau)$; each $\nu_\pp$ is Schwartz--Bruhat \cite[Lemma~A.4]{N}.

\begin{problem}\label{pr:nu-values}
  For $\tau, \mu \in F^\times$, $\nu(\tau) = m_\Phi(\tau^n / y')$ and $\hat{\nu}(\mu) = r_{\Fc \Phi}(\mu^n y')$; moreover $\hat{\nu}(0) = (b \Fc \Phi)_Y(0, y')$.
\end{problem}

For a pure tensor $s = \otimes s_\pp$ of leading terms, with $s_\pp = b \Phi_\pp^0$ for almost all $\pp$, let $T(s) := \frac{1}{n} \sum_{y \in \Delta} \prod_\pp (s_\pp)_Y(0, y)$, a finite sum.  Summing Poisson's formula for $\nu$ over $y' \in \Delta$, doing the same for the function $\sigma$ of Theorem~\ref{in-th:axis-to-marg}, and cancelling the two remaining terms at the origin by Problem~\ref{pr:isotropy}, one obtains \cite[Proposition~B.4]{N}:

\begin{theorem}\label{th:constant-terms}
  For each pure tensor $\Phi$, $C(\Phi) + T(b \Phi) = C(\Fc \Phi) + T(b \Fc \Phi)$.
\end{theorem}

Let $D(\Psi) := \Theta(\Psi) - \Theta(\Fc \Psi)$.

\begin{problem}\label{pr:translated-ct}
  If $b \Phi_\pp = 0$ for some $\pp$ and $b \Fc \Phi_{\mathfrak{q}} = 0$ for some $\mathfrak{q}$, then $\int_{\Ab / F} D(n(u) g \Phi) \, d u = 0$ for every lift $g$.
\end{problem}

Under these hypotheses, $\delta(m) := \lvert D(g \Phi) \rvert$ (for $g$ a lift of $m$) is a function on $\SL(F) \backslash \SL(\Ab)$ with vanishing constant term.  By reduction theory \cite[I.2.1]{MW} and the bounds behind the convergence of $\Theta$, it decays rapidly in the cusp.  Writing $F^2 - \{0\}$ as the $\SL(F)$-orbit of $(0, 1)$ and unfolding, $\int \delta^2$ reduces to integrals of constant terms, so $\delta = 0$ \cite[Proposition~B.6]{N}.

\begin{theorem}\label{th:poisson-vanishing}
  If $\Phi$ is a pure tensor with $b \Phi_\pp = 0$ for some $\pp$ and $b \Fc \Phi_{\mathfrak{q}} = 0$ for some $\mathfrak{q}$, then $\Theta(\Phi) = \Theta(\Fc \Phi)$.
\end{theorem}

To treat general $\Phi$, we modify one factor without changing its leading term.  Fix a finite place $w$, with residue cardinality $q$ and uniformizer $\varpi$.

\begin{problem}\label{pr:projector}
  For $\Phi_w \in \Sc_w$, $b(\Phi_w(\varpi^n \, \cdot)) = q^{n-1} b \Phi_w$ and $b \Fc(\Phi_w(\varpi^n \, \cdot)) = q^{n+1} b \Fc \Phi_w$.  Hence $Q_w \Phi_w := (q^{n+1} \Phi_w - \Phi_w(\varpi^n \, \cdot)) / (q^{n+1} - q^{n-1})$ satisfies $b Q_w \Phi_w = b \Phi_w$ and $b \Fc Q_w \Phi_w = 0$.
\end{problem}

Let $b(\Sc_\Ab) := \bigotimes'_\pp b(\Sc_\pp)$, restricted with respect to the $b \Phi_\pp^0$.  Letting $Q$ act by $Q_w$ in the $w$-factor, $D \circ Q$ vanishes on the kernel of $b$ by Theorem~\ref{th:poisson-vanishing}, so $B(b \Phi) := - D(Q \Phi)$ is well defined.  Since $D(\Fc \Psi) = - D(\Psi)$, one deduces the main theorem of \cite{N} (Theorem~9.1 and Corollary~9.2 there):

\begin{theorem}\label{th:poisson}
  There is a unique linear functional $B$ on $b(\Sc_\Ab)$ such that for all $\Phi \in \Sc_\Ab$,
  \begin{equation*}
    \sum_{v \in F^2 - \{0\}} \Phi(v) + B(b \Phi) = \sum_{v \in F^2 - \{0\}} \Fc \Phi(v) + B(b \Fc \Phi).
  \end{equation*}
  For $n = 1$, $B(b \Phi) = \prod_\pp \Phi_\pp(0)$.  If some factor of $\Phi$ lies in $\Dc$ and satisfies the vanishing condition, then $\sum_v \Phi(v) = \sum_v \Fc \Phi(v)$.
\end{theorem}

Kubota's coefficients $\qr{d}{c}_n$ appear in the sum over $v$; Gauss sums appear in its Fourier coefficients, the subject of the sequel \emph{Voronoi formulas for Gauss sums from the $n$-fold cover of $\mathrm{SL}_2$}.  Integrating Theorem~\ref{th:poisson} for $n(u) \Phi$ against $\psi(-u)$ over $\Ab / F$, the terms with $x = 0$ drop out, and those with $x = c \in F^\times$ unfold to products of local integrals.

%======================================================================
\section*{Questions}

We do not know the answers to the following questions.
\begin{enumerate}
\item What is the boundary functional $B$ of Theorem~\ref{th:poisson} for $n \geq 2$?  One expects it to be expressible through the residues of Kubota's Eisenstein series.
\item Does Theorem~\ref{th:poisson}, applied to dilates and followed by a Mellin transform, recover the continuation and functional equations of Kubota's Eisenstein series \cite{Kub69}, with the boundary terms giving the residues?  For $n = 1$ it does \cite[\S\S3--5]{God}.
\item For $F$ archimedean and $n \geq 2$, what is $\Sc / \Dc$?  Theorem~\ref{th:finite-schwartz} answers this at finite places.  For $n = 1$, $\Sc$ is the Schwartz space of $F^2$, by Estrada's theorem \cite[Theorem~4.1]{Est}.
\end{enumerate}

%======================================================================
\section*{Hints}

\subsection*{\S\ref{sec:germs}}
\begin{hints}
\item[\ref{pr:leading-transform}] Let $\mu \to 0$ in Problem~\ref{ft-pr:ft-root}.
\item[\ref{pr:finite-germs}] A bounded family of locally constant compactly supported functions on $F$ is constant on a common neighborhood of $0$.
\item[\ref{pr:germ-average}] Interchange the line integral with the integral over $a$, and use $R_X(\Phi(\mu^n \, \cdot))(v, 1) = \lvert \mu \rvert^{-n} R_X \Phi(v, \mu)$.
\end{hints}

\subsection*{\S\ref{sec:adelic}}
\begin{hints}
\item[\ref{pr:global-values}] Hilbert reciprocity and the product formula.
\item[\ref{pr:theta-S}] At $\pp \notin S$, $v$ is $\varpi_\pp^{k_\pp}$ times a primitive vector with $n \mid k_\pp$; the ideal generated by the coordinates of $v$ is an $n$th power.
\item[\ref{pr:rational-invariance}] At points with $x = 0$, use Problem~\ref{in-pr:unip-other-chart}.
\end{hints}

\subsection*{\S\ref{sec:poisson}}
\begin{hints}
\item[\ref{pr:constant-term}] The terms with $x = 0$ do not depend on $u$; unfold the others.
\item[\ref{pr:isotropy}] The transform $\beta_\pp \mapsto \lvert n \rvert_\pp^{1/2} \tilde{\beta}_\pp$ is the Fourier transform for the self-dual measure.
\item[\ref{pr:nu-values}] \eqref{in-eq:marg-to-axis}, reciprocity and the product formula; for $\hat{\nu}(0)$, the definition of the leading term.
\item[\ref{pr:translated-ct}] Problem~\ref{pr:constant-term} and Theorem~\ref{th:constant-terms}; leading terms commute with the local operators.
\item[\ref{pr:projector}] Problem~\ref{pr:leading-homog} and the dilation formula of \S\ref{ft-sec:transform}.
\end{hints}

\begin{thebibliography}{McN}

\bibitem[Est]{Est}
R.~Estrada, \emph{Vector moment problems for rapidly decreasing smooth functions of several variables}, Proc. Amer. Math. Soc. \textbf{126} (1998), no.~3, 761--768.

\bibitem[God]{God}
R.~Godement, \emph{The decomposition of {$L^2(G/\Gamma)$} for {$\Gamma = \mathrm{SL}(2, \mathbf{Z})$}}, Algebraic Groups and Discontinuous Subgroups, Proc. Sympos. Pure Math., vol.~9, Amer. Math. Soc., 1966, pp.~211--224.

\bibitem[Kub69]{Kub69}
T.~Kubota, \emph{On automorphic functions and the reciprocity law in a number field}, Lectures in Mathematics, Department of Mathematics, Kyoto University, No.~2, Kinokuniya, 1969.

\bibitem[MW]{MW}
C.~M{\oe}glin and J.-L. Waldspurger, \emph{Spectral decomposition and {E}isenstein series}, Cambridge Tracts in Mathematics, vol.~113, Cambridge University Press, 1995.

\bibitem[N]{N}
P.~D. Nelson, \emph{A {S}chwartz space for the $n$-fold cover of $\mathrm{SL}_2$}, preprint, 2026.

\bibitem[Neu]{Neu}
J.~Neukirch, \emph{Algebraic number theory}, Grundlehren der mathematischen Wissenschaften, vol.~322, Springer, 1999.

\bibitem[Tate]{Tate}
J.~T. Tate, \emph{Fourier analysis in number fields, and {H}ecke's zeta-functions}, Algebraic Number Theory (Proc. Instructional Conf., Brighton, 1965), Thompson, 1967, pp.~305--347.

\bibitem[Weil]{Weil}
A.~Weil, \emph{Basic number theory}, third ed., Grundlehren der mathematischen Wissenschaften, vol.~144, Springer, 1974.

\end{thebibliography}

\end{document}
